% 2013 年安徽卷文科第 3 题：输出 s 的程序框图（算法流程图）。
% 来源 PDF：普通高考/2013/2013安徽文.pdf，第 1 页（印面 785）左栏。
% 流程：开始 → s=0,n=2 → n<8? 是：s=s+1/n，n=n+2，回到判断；否：输出 s → 结束。
\begin{tikzpicture}[semithick, >=Stealth, font=\small,
  block/.style={draw, minimum width=2.9cm, minimum height=0.62cm, align=center},
  dec/.style={draw, diamond, aspect=2.5, inner sep=1pt, align=center},
  term/.style={draw, rounded corners=6pt, minimum width=1.4cm, minimum height=0.52cm},
  io/.style={draw, trapezium, trapezium left angle=70, trapezium right angle=110,
             minimum width=2.3cm, minimum height=0.62cm},
  arr/.style={->}]
  \node (start) [term] at (0,0) {开始};
  \node (init)  [block] at (0,-1.05) {$s=0,\ n=2$};
  \node (dec)   [dec]  at (0,-2.55) {$n<8$?};
  \node (body)  [block] at (0,-4.0) {$s=s+\dfrac1n$};
  \node (inc)   [block] at (0,-5.3) {$n=n+2$};
  \node (out)   [io]   at (3.2,-4.0) {输出 $s$};
  \node (end)   [term] at (3.2,-5.3) {结束};
  \draw[arr] (start) -- (init);
  \draw[arr] (init) -- (dec);
  \draw[arr] (dec) -- node[left, xshift=-1pt] {是} (body);
  \draw[arr] (body) -- (inc);
  \draw[arr] (dec.east) -- node[above] {否} (3.2,-2.55) -- (out.north);
  \draw[arr] (out) -- (end);
  % 是分支的回路：n=n+2 后向左、向上，再向右汇入判断框上方的流程线
  \draw[arr] (inc.west) -- (-2.0,-5.3) -- (-2.0,-1.8) -- (0,-1.8);
\end{tikzpicture}
